\documentclass[aspectratio=169,10pt]{beamer}
% ---------------------------------------------------------------
% R3.07 SQL dans un langage de programmation -- preambule commun des diapositives
% Charte derivee de R5.A.08 Qualite de developpement (metropolis, dark tech)
% Palette IUT partagee (violet du logo + accent or)
% Utilisation :
%   \documentclass[aspectratio=169,10pt]{beamer}
%   % ---------------------------------------------------------------
% R3.07 SQL dans un langage de programmation -- preambule commun des diapositives
% Charte derivee de R5.A.08 Qualite de developpement (metropolis, dark tech)
% Palette IUT partagee (violet du logo + accent or)
% Utilisation :
%   \documentclass[aspectratio=169,10pt]{beamer}
%   % ---------------------------------------------------------------
% R3.07 SQL dans un langage de programmation -- preambule commun des diapositives
% Charte derivee de R5.A.08 Qualite de developpement (metropolis, dark tech)
% Palette IUT partagee (violet du logo + accent or)
% Utilisation :
%   \documentclass[aspectratio=169,10pt]{beamer}
%   \input{../../latex/common/r307-slides-preamble.tex}
%   \graphicspath{{../../latex/common/assets/}{img/}{figs/}}
% ---------------------------------------------------------------
\usetheme{metropolis}
\usepackage{appendixnumberbeamer}

\usepackage[T1]{fontenc}
\usepackage[utf8]{inputenc}
\usepackage[french]{babel}
\usepackage[scale=0.94]{tgheros}
\renewcommand{\familydefault}{\sfdefault}
\usepackage{amsmath,amssymb,amsfonts}
\usepackage{graphicx}
\usepackage{booktabs}
\usepackage{multirow}
\usepackage{tikz}
\usetikzlibrary{shapes.geometric, shapes, arrows, arrows.meta, positioning, calc, fit}
\usepackage{xcolor}
\usepackage{listings}
\usepackage[skins,listings]{tcolorbox}
\usepackage{qrcode}
\usepackage{hyperref}

% ---------------------------------------------------------------
% PALETTE (dark-mode tech: near-black canvas, violet (logo IUT) + gold accents)
% ---------------------------------------------------------------
\definecolor{qabg}{HTML}{0E1016}      % canvas background
\definecolor{qapanel}{HTML}{181B24}   % card / block panel
\definecolor{qapanelalt}{HTML}{20242F}% secondary panel tint
\definecolor{qaline}{HTML}{2B303D}    % hairline rules / table rules
\definecolor{qatext}{HTML}{ECEEF2}    % primary text (off-white)
\definecolor{qamuted}{HTML}{8B93A6}   % secondary / caption text
\definecolor{qaaccent}{HTML}{D93ED7}  % primary accent -- violet du logo IUT, eclairci pour fond sombre
\definecolor{qaaccent2}{HTML}{F2B84B} % secondary accent -- or/ambre
\definecolor{qacodekw}{HTML}{C792EA}  % code : mots-cles
\definecolor{qacodestr}{HTML}{C3E88D} % code : chaines
\definecolor{qacodecom}{HTML}{7F8AA3} % code : commentaires
\definecolor{qacodeann}{HTML}{F2B84B} % code : annotations / javadoc

\setbeamercolor{background canvas}{bg=qabg}
\setbeamercolor{normal text}{fg=qatext}
\setbeamercolor{palette primary}{bg=qabg, fg=qamuted}
\setbeamercolor{frametitle}{bg=qabg, fg=qatext}
\setbeamercolor{progress bar}{fg=qaaccent, bg=qaline}
\setbeamercolor{alerted text}{fg=qaaccent2}
\setbeamercolor{block title}{bg=qapanel, fg=qaaccent}
\setbeamercolor{block body}{bg=qapanel, fg=qatext}
\setbeamercolor{block title alerted}{bg=qapanel, fg=qaaccent2}
\setbeamercolor{block body alerted}{bg=qapanel, fg=qatext}
\setbeamercolor{block title example}{bg=qapanelalt, fg=qaaccent}
\setbeamercolor{block body example}{bg=qapanelalt, fg=qatext}
\setbeamercolor{example text}{fg=qaaccent}
\setbeamercolor{itemize item}{fg=qaaccent}
\setbeamercolor{itemize subitem}{fg=qaaccent2}
\setbeamercolor{enumerate item}{fg=qaaccent}
\setbeamercolor{enumerate subitem}{fg=qaaccent2}
\setbeamercolor{caption}{fg=qamuted}
\setbeamercolor{page number in head/foot}{fg=qamuted, bg=qabg}
\setbeamercolor{title}{fg=qatext}
\setbeamercolor{subtitle}{fg=qamuted}
\setbeamercolor{author}{fg=qatext}
\setbeamercolor{institute}{fg=qamuted}
\setbeamercolor{date}{fg=qamuted}
\setbeamercolor{section in toc}{fg=qatext}
\setbeamercolor{subsection in toc}{fg=qamuted}

\setbeamertemplate{blocks}[rounded][shadow=false]
\setbeamertemplate{itemize item}{\raisebox{0.15em}{\scriptsize\textcolor{qaaccent}{\ensuremath{\blacktriangleright}}}}
\setbeamertemplate{itemize subitem}{\raisebox{0.1em}{\tiny\textcolor{qaaccent2}{\ensuremath{\bullet}}}}

\setbeamerfont{frametitle}{series=\bfseries, size=\Large}
\setbeamerfont{title}{series=\bfseries, size=\Large}
\setbeamerfont{subtitle}{size=\normalsize}
\setbeamerfont{institute}{size=\scriptsize}
\setbeamerfont{date}{size=\scriptsize}
\setbeamerfont{author}{size=\small}

\metroset{sectionpage=none, subsectionpage=none, numbering=fraction, progressbar=frametitle}
\setbeamertemplate{section in toc}[sections numbered]

% ---------------------------------------------------------------
% LISTINGS -- theme sombre (lstlisting reste utilisable tel quel)
% \lstcode : police + couleur de base, a reutiliser dans les
%            surcharges locales (ex. basicstyle=\lstcode\scriptsize)
% ---------------------------------------------------------------
\newcommand{\lstcode}{\ttfamily\color{qatext}}

\lstdefinelanguage{TypeScript}{
  keywords={typeof, new, true, false, catch, function, return, null, switch, var,
    if, in, while, do, else, case, break, const, let, interface, class, extends,
    implements, enum, export, import, this, super, constructor, public, private,
    protected, static, async, await, throw, throws, try, finally, void},
  ndkeywords={boolean, any, number, string, int, long, double},
  sensitive=true,
  comment=[l]{//},
  morecomment=[s]{/*}{*/},
  morestring=[b]',
  morestring=[b]"
}

\lstset{
  language=Java,
  basicstyle=\lstcode\footnotesize,
  keywordstyle=\color{qacodekw}\bfseries,
  ndkeywordstyle=\color{qaaccent2}\bfseries,
  identifierstyle=\color{qatext},
  stringstyle=\color{qacodestr},
  commentstyle=\color{qacodecom}\itshape,
  morecomment=[s][\color{qacodeann}]{/**}{*/},
  backgroundcolor=\color{qapanelalt},
  rulecolor=\color{qaline},
  frame=single,
  framesep=3pt,
  xleftmargin=4pt,
  xrightmargin=4pt,
  keepspaces=true,
  columns=fullflexible,
  showstringspaces=false,
  breaklines=true,
  tabsize=2,
  escapeinside={(*@}{@*)},
}

% ---------------------------------------------------------------
% CARD-STYLE BOXES
% ---------------------------------------------------------------
\tcbset{
  figcard/.style={colback=white, colframe=white, boxrule=0pt, arc=4pt,
    left=3pt, right=3pt, top=3pt, bottom=3pt, boxsep=0pt,
    halign=center, valign=center},
  datacard/.style={colback=qapanel, colframe=qaline, boxrule=0.4pt, arc=3pt,
    left=8pt, right=8pt, top=6pt, bottom=6pt, coltext=qatext,
    before upper={\color{qatext}}, halign=center},
  codebox/.style={colback=qapanelalt, colframe=qaline, boxrule=0.4pt, arc=3pt,
    left=8pt, right=8pt, top=5pt, bottom=5pt, listing only,
    listing options={basicstyle=\ttfamily\small\color{qaaccent},
      breaklines=true, showstringspaces=false, columns=fullflexible,
      keepspaces=true, frame=none, backgroundcolor={}}}
}
\newtcolorbox{figcard}[1][]{figcard,#1}
\newtcolorbox{datacard}{datacard}
\newtcblisting{codebox}{codebox}

% ---------------------------------------------------------------
% HANDY MACROS
% ---------------------------------------------------------------
\newcommand{\hl}[1]{\alert{#1}}
\newcommand{\source}[1]{{\scriptsize\color{qamuted}#1}}

% Sommaire automatique en debut de section
\AtBeginSection[]{%
  \begin{frame}<beamer>{Plan du cours}
    \tableofcontents[currentsection]
  \end{frame}}

% Page de titre : \rtitleslide{numero de chapitre}{titre du chapitre}
% Les auteurs sont repris de \rauthors (modifiable avant \begin{document})
\newcommand{\rauthors}{Selena Lamari et Chouki Tibermacine}
\newcommand{\rtitleslide}[2]{%
\begin{frame}[plain,noframenumbering]
  \begin{center}
    \vspace{2mm}
    {\color{qaaccent}\rule{2.1cm}{1.2pt}}\\[6mm]
    {\Large\bfseries\color{qatext} R3.07}\\[1mm]
    {\Large\bfseries\color{qatext} SQL dans un langage de programmation}\\[5mm]
    {\large\color{qamuted} Chapitre #1}\\[1mm]
    {\large\color{qaaccent} #2}\\[6mm]
    \includegraphics[height=2.1cm]{iut_logo_transparent.png}

    \vspace{4mm}

    {\small\color{qatext} \rauthors}\\[1mm]
    {\scriptsize\color{qamuted} BUT 2 Informatique --- IUT de Vannes --- Universit\'e Bretagne Sud}
  \end{center}
\end{frame}
}

% ---- Diapo de fin « Questions ? » (même modèle que R5.A.08) ----
\newcommand{\rcontact}{Chouki.Tibermacine@univ-ubs.fr}
\newcommand{\rquestionslide}{%
\begin{frame}{Questions ?}
  \begin{center}
    \vspace{2mm}
    \begin{tikzpicture}
      \node[regular polygon, regular polygon sides=6, draw=qaline, thick,
        minimum size=4.6cm, rotate=90] (hexin) {};
      \node[regular polygon, regular polygon sides=6, draw=qaaccent, dashed,
        thick, minimum size=6.2cm, rotate=90] (hexout) {};
      \node[font=\fontsize{92}{92}\selectfont\bfseries, color=qaaccent2] at (hexin) {?};
    \end{tikzpicture}\\[5mm]
    {\small\color{qamuted} \rcontact}
  \end{center}
\end{frame}}

%   \graphicspath{{../../latex/common/assets/}{img/}{figs/}}
% ---------------------------------------------------------------
\usetheme{metropolis}
\usepackage{appendixnumberbeamer}

\usepackage[T1]{fontenc}
\usepackage[utf8]{inputenc}
\usepackage[french]{babel}
\usepackage[scale=0.94]{tgheros}
\renewcommand{\familydefault}{\sfdefault}
\usepackage{amsmath,amssymb,amsfonts}
\usepackage{graphicx}
\usepackage{booktabs}
\usepackage{multirow}
\usepackage{tikz}
\usetikzlibrary{shapes.geometric, shapes, arrows, arrows.meta, positioning, calc, fit}
\usepackage{xcolor}
\usepackage{listings}
\usepackage[skins,listings]{tcolorbox}
\usepackage{qrcode}
\usepackage{hyperref}

% ---------------------------------------------------------------
% PALETTE (dark-mode tech: near-black canvas, violet (logo IUT) + gold accents)
% ---------------------------------------------------------------
\definecolor{qabg}{HTML}{0E1016}      % canvas background
\definecolor{qapanel}{HTML}{181B24}   % card / block panel
\definecolor{qapanelalt}{HTML}{20242F}% secondary panel tint
\definecolor{qaline}{HTML}{2B303D}    % hairline rules / table rules
\definecolor{qatext}{HTML}{ECEEF2}    % primary text (off-white)
\definecolor{qamuted}{HTML}{8B93A6}   % secondary / caption text
\definecolor{qaaccent}{HTML}{D93ED7}  % primary accent -- violet du logo IUT, eclairci pour fond sombre
\definecolor{qaaccent2}{HTML}{F2B84B} % secondary accent -- or/ambre
\definecolor{qacodekw}{HTML}{C792EA}  % code : mots-cles
\definecolor{qacodestr}{HTML}{C3E88D} % code : chaines
\definecolor{qacodecom}{HTML}{7F8AA3} % code : commentaires
\definecolor{qacodeann}{HTML}{F2B84B} % code : annotations / javadoc

\setbeamercolor{background canvas}{bg=qabg}
\setbeamercolor{normal text}{fg=qatext}
\setbeamercolor{palette primary}{bg=qabg, fg=qamuted}
\setbeamercolor{frametitle}{bg=qabg, fg=qatext}
\setbeamercolor{progress bar}{fg=qaaccent, bg=qaline}
\setbeamercolor{alerted text}{fg=qaaccent2}
\setbeamercolor{block title}{bg=qapanel, fg=qaaccent}
\setbeamercolor{block body}{bg=qapanel, fg=qatext}
\setbeamercolor{block title alerted}{bg=qapanel, fg=qaaccent2}
\setbeamercolor{block body alerted}{bg=qapanel, fg=qatext}
\setbeamercolor{block title example}{bg=qapanelalt, fg=qaaccent}
\setbeamercolor{block body example}{bg=qapanelalt, fg=qatext}
\setbeamercolor{example text}{fg=qaaccent}
\setbeamercolor{itemize item}{fg=qaaccent}
\setbeamercolor{itemize subitem}{fg=qaaccent2}
\setbeamercolor{enumerate item}{fg=qaaccent}
\setbeamercolor{enumerate subitem}{fg=qaaccent2}
\setbeamercolor{caption}{fg=qamuted}
\setbeamercolor{page number in head/foot}{fg=qamuted, bg=qabg}
\setbeamercolor{title}{fg=qatext}
\setbeamercolor{subtitle}{fg=qamuted}
\setbeamercolor{author}{fg=qatext}
\setbeamercolor{institute}{fg=qamuted}
\setbeamercolor{date}{fg=qamuted}
\setbeamercolor{section in toc}{fg=qatext}
\setbeamercolor{subsection in toc}{fg=qamuted}

\setbeamertemplate{blocks}[rounded][shadow=false]
\setbeamertemplate{itemize item}{\raisebox{0.15em}{\scriptsize\textcolor{qaaccent}{\ensuremath{\blacktriangleright}}}}
\setbeamertemplate{itemize subitem}{\raisebox{0.1em}{\tiny\textcolor{qaaccent2}{\ensuremath{\bullet}}}}

\setbeamerfont{frametitle}{series=\bfseries, size=\Large}
\setbeamerfont{title}{series=\bfseries, size=\Large}
\setbeamerfont{subtitle}{size=\normalsize}
\setbeamerfont{institute}{size=\scriptsize}
\setbeamerfont{date}{size=\scriptsize}
\setbeamerfont{author}{size=\small}

\metroset{sectionpage=none, subsectionpage=none, numbering=fraction, progressbar=frametitle}
\setbeamertemplate{section in toc}[sections numbered]

% ---------------------------------------------------------------
% LISTINGS -- theme sombre (lstlisting reste utilisable tel quel)
% \lstcode : police + couleur de base, a reutiliser dans les
%            surcharges locales (ex. basicstyle=\lstcode\scriptsize)
% ---------------------------------------------------------------
\newcommand{\lstcode}{\ttfamily\color{qatext}}

\lstdefinelanguage{TypeScript}{
  keywords={typeof, new, true, false, catch, function, return, null, switch, var,
    if, in, while, do, else, case, break, const, let, interface, class, extends,
    implements, enum, export, import, this, super, constructor, public, private,
    protected, static, async, await, throw, throws, try, finally, void},
  ndkeywords={boolean, any, number, string, int, long, double},
  sensitive=true,
  comment=[l]{//},
  morecomment=[s]{/*}{*/},
  morestring=[b]',
  morestring=[b]"
}

\lstset{
  language=Java,
  basicstyle=\lstcode\footnotesize,
  keywordstyle=\color{qacodekw}\bfseries,
  ndkeywordstyle=\color{qaaccent2}\bfseries,
  identifierstyle=\color{qatext},
  stringstyle=\color{qacodestr},
  commentstyle=\color{qacodecom}\itshape,
  morecomment=[s][\color{qacodeann}]{/**}{*/},
  backgroundcolor=\color{qapanelalt},
  rulecolor=\color{qaline},
  frame=single,
  framesep=3pt,
  xleftmargin=4pt,
  xrightmargin=4pt,
  keepspaces=true,
  columns=fullflexible,
  showstringspaces=false,
  breaklines=true,
  tabsize=2,
  escapeinside={(*@}{@*)},
}

% ---------------------------------------------------------------
% CARD-STYLE BOXES
% ---------------------------------------------------------------
\tcbset{
  figcard/.style={colback=white, colframe=white, boxrule=0pt, arc=4pt,
    left=3pt, right=3pt, top=3pt, bottom=3pt, boxsep=0pt,
    halign=center, valign=center},
  datacard/.style={colback=qapanel, colframe=qaline, boxrule=0.4pt, arc=3pt,
    left=8pt, right=8pt, top=6pt, bottom=6pt, coltext=qatext,
    before upper={\color{qatext}}, halign=center},
  codebox/.style={colback=qapanelalt, colframe=qaline, boxrule=0.4pt, arc=3pt,
    left=8pt, right=8pt, top=5pt, bottom=5pt, listing only,
    listing options={basicstyle=\ttfamily\small\color{qaaccent},
      breaklines=true, showstringspaces=false, columns=fullflexible,
      keepspaces=true, frame=none, backgroundcolor={}}}
}
\newtcolorbox{figcard}[1][]{figcard,#1}
\newtcolorbox{datacard}{datacard}
\newtcblisting{codebox}{codebox}

% ---------------------------------------------------------------
% HANDY MACROS
% ---------------------------------------------------------------
\newcommand{\hl}[1]{\alert{#1}}
\newcommand{\source}[1]{{\scriptsize\color{qamuted}#1}}

% Sommaire automatique en debut de section
\AtBeginSection[]{%
  \begin{frame}<beamer>{Plan du cours}
    \tableofcontents[currentsection]
  \end{frame}}

% Page de titre : \rtitleslide{numero de chapitre}{titre du chapitre}
% Les auteurs sont repris de \rauthors (modifiable avant \begin{document})
\newcommand{\rauthors}{Selena Lamari et Chouki Tibermacine}
\newcommand{\rtitleslide}[2]{%
\begin{frame}[plain,noframenumbering]
  \begin{center}
    \vspace{2mm}
    {\color{qaaccent}\rule{2.1cm}{1.2pt}}\\[6mm]
    {\Large\bfseries\color{qatext} R3.07}\\[1mm]
    {\Large\bfseries\color{qatext} SQL dans un langage de programmation}\\[5mm]
    {\large\color{qamuted} Chapitre #1}\\[1mm]
    {\large\color{qaaccent} #2}\\[6mm]
    \includegraphics[height=2.1cm]{iut_logo_transparent.png}

    \vspace{4mm}

    {\small\color{qatext} \rauthors}\\[1mm]
    {\scriptsize\color{qamuted} BUT 2 Informatique --- IUT de Vannes --- Universit\'e Bretagne Sud}
  \end{center}
\end{frame}
}

% ---- Diapo de fin « Questions ? » (même modèle que R5.A.08) ----
\newcommand{\rcontact}{Chouki.Tibermacine@univ-ubs.fr}
\newcommand{\rquestionslide}{%
\begin{frame}{Questions ?}
  \begin{center}
    \vspace{2mm}
    \begin{tikzpicture}
      \node[regular polygon, regular polygon sides=6, draw=qaline, thick,
        minimum size=4.6cm, rotate=90] (hexin) {};
      \node[regular polygon, regular polygon sides=6, draw=qaaccent, dashed,
        thick, minimum size=6.2cm, rotate=90] (hexout) {};
      \node[font=\fontsize{92}{92}\selectfont\bfseries, color=qaaccent2] at (hexin) {?};
    \end{tikzpicture}\\[5mm]
    {\small\color{qamuted} \rcontact}
  \end{center}
\end{frame}}

%   \graphicspath{{../../latex/common/assets/}{img/}{figs/}}
% ---------------------------------------------------------------
\usetheme{metropolis}
\usepackage{appendixnumberbeamer}

\usepackage[T1]{fontenc}
\usepackage[utf8]{inputenc}
\usepackage[french]{babel}
\usepackage[scale=0.94]{tgheros}
\renewcommand{\familydefault}{\sfdefault}
\usepackage{amsmath,amssymb,amsfonts}
\usepackage{graphicx}
\usepackage{booktabs}
\usepackage{multirow}
\usepackage{tikz}
\usetikzlibrary{shapes.geometric, shapes, arrows, arrows.meta, positioning, calc, fit}
\usepackage{xcolor}
\usepackage{listings}
\usepackage[skins,listings]{tcolorbox}
\usepackage{qrcode}
\usepackage{hyperref}

% ---------------------------------------------------------------
% PALETTE (dark-mode tech: near-black canvas, violet (logo IUT) + gold accents)
% ---------------------------------------------------------------
\definecolor{qabg}{HTML}{0E1016}      % canvas background
\definecolor{qapanel}{HTML}{181B24}   % card / block panel
\definecolor{qapanelalt}{HTML}{20242F}% secondary panel tint
\definecolor{qaline}{HTML}{2B303D}    % hairline rules / table rules
\definecolor{qatext}{HTML}{ECEEF2}    % primary text (off-white)
\definecolor{qamuted}{HTML}{8B93A6}   % secondary / caption text
\definecolor{qaaccent}{HTML}{D93ED7}  % primary accent -- violet du logo IUT, eclairci pour fond sombre
\definecolor{qaaccent2}{HTML}{F2B84B} % secondary accent -- or/ambre
\definecolor{qacodekw}{HTML}{C792EA}  % code : mots-cles
\definecolor{qacodestr}{HTML}{C3E88D} % code : chaines
\definecolor{qacodecom}{HTML}{7F8AA3} % code : commentaires
\definecolor{qacodeann}{HTML}{F2B84B} % code : annotations / javadoc

\setbeamercolor{background canvas}{bg=qabg}
\setbeamercolor{normal text}{fg=qatext}
\setbeamercolor{palette primary}{bg=qabg, fg=qamuted}
\setbeamercolor{frametitle}{bg=qabg, fg=qatext}
\setbeamercolor{progress bar}{fg=qaaccent, bg=qaline}
\setbeamercolor{alerted text}{fg=qaaccent2}
\setbeamercolor{block title}{bg=qapanel, fg=qaaccent}
\setbeamercolor{block body}{bg=qapanel, fg=qatext}
\setbeamercolor{block title alerted}{bg=qapanel, fg=qaaccent2}
\setbeamercolor{block body alerted}{bg=qapanel, fg=qatext}
\setbeamercolor{block title example}{bg=qapanelalt, fg=qaaccent}
\setbeamercolor{block body example}{bg=qapanelalt, fg=qatext}
\setbeamercolor{example text}{fg=qaaccent}
\setbeamercolor{itemize item}{fg=qaaccent}
\setbeamercolor{itemize subitem}{fg=qaaccent2}
\setbeamercolor{enumerate item}{fg=qaaccent}
\setbeamercolor{enumerate subitem}{fg=qaaccent2}
\setbeamercolor{caption}{fg=qamuted}
\setbeamercolor{page number in head/foot}{fg=qamuted, bg=qabg}
\setbeamercolor{title}{fg=qatext}
\setbeamercolor{subtitle}{fg=qamuted}
\setbeamercolor{author}{fg=qatext}
\setbeamercolor{institute}{fg=qamuted}
\setbeamercolor{date}{fg=qamuted}
\setbeamercolor{section in toc}{fg=qatext}
\setbeamercolor{subsection in toc}{fg=qamuted}

\setbeamertemplate{blocks}[rounded][shadow=false]
\setbeamertemplate{itemize item}{\raisebox{0.15em}{\scriptsize\textcolor{qaaccent}{\ensuremath{\blacktriangleright}}}}
\setbeamertemplate{itemize subitem}{\raisebox{0.1em}{\tiny\textcolor{qaaccent2}{\ensuremath{\bullet}}}}

\setbeamerfont{frametitle}{series=\bfseries, size=\Large}
\setbeamerfont{title}{series=\bfseries, size=\Large}
\setbeamerfont{subtitle}{size=\normalsize}
\setbeamerfont{institute}{size=\scriptsize}
\setbeamerfont{date}{size=\scriptsize}
\setbeamerfont{author}{size=\small}

\metroset{sectionpage=none, subsectionpage=none, numbering=fraction, progressbar=frametitle}
\setbeamertemplate{section in toc}[sections numbered]

% ---------------------------------------------------------------
% LISTINGS -- theme sombre (lstlisting reste utilisable tel quel)
% \lstcode : police + couleur de base, a reutiliser dans les
%            surcharges locales (ex. basicstyle=\lstcode\scriptsize)
% ---------------------------------------------------------------
\newcommand{\lstcode}{\ttfamily\color{qatext}}

\lstdefinelanguage{TypeScript}{
  keywords={typeof, new, true, false, catch, function, return, null, switch, var,
    if, in, while, do, else, case, break, const, let, interface, class, extends,
    implements, enum, export, import, this, super, constructor, public, private,
    protected, static, async, await, throw, throws, try, finally, void},
  ndkeywords={boolean, any, number, string, int, long, double},
  sensitive=true,
  comment=[l]{//},
  morecomment=[s]{/*}{*/},
  morestring=[b]',
  morestring=[b]"
}

\lstset{
  language=Java,
  basicstyle=\lstcode\footnotesize,
  keywordstyle=\color{qacodekw}\bfseries,
  ndkeywordstyle=\color{qaaccent2}\bfseries,
  identifierstyle=\color{qatext},
  stringstyle=\color{qacodestr},
  commentstyle=\color{qacodecom}\itshape,
  morecomment=[s][\color{qacodeann}]{/**}{*/},
  backgroundcolor=\color{qapanelalt},
  rulecolor=\color{qaline},
  frame=single,
  framesep=3pt,
  xleftmargin=4pt,
  xrightmargin=4pt,
  keepspaces=true,
  columns=fullflexible,
  showstringspaces=false,
  breaklines=true,
  tabsize=2,
  escapeinside={(*@}{@*)},
}

% ---------------------------------------------------------------
% CARD-STYLE BOXES
% ---------------------------------------------------------------
\tcbset{
  figcard/.style={colback=white, colframe=white, boxrule=0pt, arc=4pt,
    left=3pt, right=3pt, top=3pt, bottom=3pt, boxsep=0pt,
    halign=center, valign=center},
  datacard/.style={colback=qapanel, colframe=qaline, boxrule=0.4pt, arc=3pt,
    left=8pt, right=8pt, top=6pt, bottom=6pt, coltext=qatext,
    before upper={\color{qatext}}, halign=center},
  codebox/.style={colback=qapanelalt, colframe=qaline, boxrule=0.4pt, arc=3pt,
    left=8pt, right=8pt, top=5pt, bottom=5pt, listing only,
    listing options={basicstyle=\ttfamily\small\color{qaaccent},
      breaklines=true, showstringspaces=false, columns=fullflexible,
      keepspaces=true, frame=none, backgroundcolor={}}}
}
\newtcolorbox{figcard}[1][]{figcard,#1}
\newtcolorbox{datacard}{datacard}
\newtcblisting{codebox}{codebox}

% ---------------------------------------------------------------
% HANDY MACROS
% ---------------------------------------------------------------
\newcommand{\hl}[1]{\alert{#1}}
\newcommand{\source}[1]{{\scriptsize\color{qamuted}#1}}

% Sommaire automatique en debut de section
\AtBeginSection[]{%
  \begin{frame}<beamer>{Plan du cours}
    \tableofcontents[currentsection]
  \end{frame}}

% Page de titre : \rtitleslide{numero de chapitre}{titre du chapitre}
% Les auteurs sont repris de \rauthors (modifiable avant \begin{document})
\newcommand{\rauthors}{Selena Lamari et Chouki Tibermacine}
\newcommand{\rtitleslide}[2]{%
\begin{frame}[plain,noframenumbering]
  \begin{center}
    \vspace{2mm}
    {\color{qaaccent}\rule{2.1cm}{1.2pt}}\\[6mm]
    {\Large\bfseries\color{qatext} R3.07}\\[1mm]
    {\Large\bfseries\color{qatext} SQL dans un langage de programmation}\\[5mm]
    {\large\color{qamuted} Chapitre #1}\\[1mm]
    {\large\color{qaaccent} #2}\\[6mm]
    \includegraphics[height=2.1cm]{iut_logo_transparent.png}

    \vspace{4mm}

    {\small\color{qatext} \rauthors}\\[1mm]
    {\scriptsize\color{qamuted} BUT 2 Informatique --- IUT de Vannes --- Universit\'e Bretagne Sud}
  \end{center}
\end{frame}
}

% ---- Diapo de fin « Questions ? » (même modèle que R5.A.08) ----
\newcommand{\rcontact}{Chouki.Tibermacine@univ-ubs.fr}
\newcommand{\rquestionslide}{%
\begin{frame}{Questions ?}
  \begin{center}
    \vspace{2mm}
    \begin{tikzpicture}
      \node[regular polygon, regular polygon sides=6, draw=qaline, thick,
        minimum size=4.6cm, rotate=90] (hexin) {};
      \node[regular polygon, regular polygon sides=6, draw=qaaccent, dashed,
        thick, minimum size=6.2cm, rotate=90] (hexout) {};
      \node[font=\fontsize{92}{92}\selectfont\bfseries, color=qaaccent2] at (hexin) {?};
    \end{tikzpicture}\\[5mm]
    {\small\color{qamuted} \rcontact}
  \end{center}
\end{frame}}

\graphicspath{{../../latex/common/assets/}}

\title{Transactions et concurrence d'acc\`es}
\author[C. Tibermacine]{Selena Lamari et Chouki Tibermacine}
\date{}

\tikzset{
  tx/.style={rectangle, rounded corners, draw=qaaccent, thick, fill=qapanel, text=qatext,
    minimum width=3.6cm, minimum height=0.55cm, align=center, font=\scriptsize\ttfamily},
  txb/.style={tx, draw=qaaccent2},
  db/.style={rectangle, rounded corners, draw=qaline, fill=qapanelalt, text=qatext,
    minimum width=2.4cm, minimum height=0.55cm, align=center, font=\scriptsize},
  timeline/.style={-Latex, thick, qamuted}
}

\begin{document}

\rtitleslide{2}{Transactions et concurrence d'acc\`es}

% ================================================================
\section{Transactions : rappels}

\begin{frame}{Qu'est-ce qu'une transaction ?}
  \begin{block}{D\'efinition}
    Une \textbf{transaction} est une suite d'op\'erations sur la base qui forme une \textbf{unit\'e logique de travail} :
    soit tout est appliqu\'e (\texttt{COMMIT}), soit rien (\texttt{ROLLBACK}).
  \end{block}
  \begin{columns}[T]
    \begin{column}{0.48\textwidth}
      \begin{exampleblock}{Propri\'et\'es ACID}
        \begin{itemize}
          \item \textbf{A}tomicit\'e : tout ou rien
          \item \textbf{C}oh\'erence : les contraintes restent vraies
          \item \textbf{I}solation : les transactions concurrentes ne se g\^enent pas (\alert{\`a quel point ?})
          \item \textbf{D}urabilit\'e : un \texttt{COMMIT} survit \`a une panne
        \end{itemize}
      \end{exampleblock}
    \end{column}
    \begin{column}{0.48\textwidth}
      \begin{exampleblock}{Exemple : commande en ligne}
        \begin{enumerate}
          \item retirer les articles du stock
          \item cr\'eer la commande
          \item journaliser l'op\'eration
        \end{enumerate}
        Une panne entre 1 et 2 ne doit pas faire \og dispara\^itre \fg{} des articles du stock.
      \end{exampleblock}
    \end{column}
  \end{columns}
\end{frame}

\begin{frame}[fragile]{Transactions en JDBC : le sch\'ema \`a retenir}
\begin{lstlisting}[language=Java]
try (Connection c = dataSource.getConnection()) {
  c.setAutoCommit(false);              // par defaut : autocommit apres CHAQUE requete
  try {
    retirerStock(c, produit, quantite);
    creerCommande(c, client, produit, quantite);
    c.commit();                        // tout est valide
  } catch (SQLException e) {
    c.rollback();                      // rien n'est applique
    throw e;
  }
}                                      // close() rend la connexion au pool
\end{lstlisting}
  \begin{itemize}
    \item Toutes les requ\^etes d'une transaction utilisent la \alert{m\^eme} \texttt{Connection}
    \item Points de reprise partiels possibles : \texttt{Savepoint sp = c.setSavepoint(); \ldots{} c.rollback(sp);}
  \end{itemize}
\end{frame}

\begin{frame}[fragile]{D'o\`u vient la \texttt{Connection} ?}
  \begin{columns}[T]
    \begin{column}{0.48\textwidth}
      \textbf{En TP : \texttt{DriverManager}}
\begin{lstlisting}[language=Java]
Connection c = DriverManager
  .getConnection(url, user, pwd);
// nouvelle connexion physique
\end{lstlisting}
    \end{column}
    \begin{column}{0.48\textwidth}
      \textbf{En production : un pool}
\begin{lstlisting}[language=Java]
Connection c = dataSource
  .getConnection();   // empruntee
c.close();            // rendue au pool
\end{lstlisting}
    \end{column}
  \end{columns}
  \vspace{2mm}
  \begin{itemize}
    \item Ouvrir une connexion co\^ute cher : un \textbf{pool} (ex. HikariCP) en garde quelques-unes pr\^etes
    \item Le code transactionnel (\texttt{setAutoCommit}, \texttt{commit}, \texttt{rollback}) est \textbf{identique}
    \item Les slides de ce chapitre \'ecrivent \texttt{dataSource.getConnection()} : en TP, utilisez votre \texttt{getConnection()}
  \end{itemize}
\end{frame}

% ================================================================
\section{Concurrence d'acc\`es}

\begin{frame}{Pourquoi la concurrence est in\'evitable}
  \begin{itemize}
    \item Une application web ou mobile sert \textbf{plusieurs utilisateurs en m\^eme temps}
    \item Chaque requ\^ete HTTP est trait\'ee par un \textbf{thread} diff\'erent, avec sa propre connexion (pool)
    \item Plusieurs transactions lisent et \'ecrivent donc les \textbf{m\^emes lignes} au m\^eme moment
  \end{itemize}
  \vspace{2mm}
  \begin{alertblock}{Question}
    Deux clients commandent le m\^eme article (stock : 10) au m\^eme instant : $-2$ et $-3$.
    Quel est le stock final ? \only<2>{\alert{Pas forc\'ement 5\ldots}}
  \end{alertblock}
\end{frame}

\begin{frame}{L'anomalie la plus courante : la mise \`a jour perdue}
  \begin{center}
  \begin{tikzpicture}[node distance=0.25cm]
    \node[font=\small\bfseries, text=qaaccent] (h1) at (0,0) {Transaction T1 (client 1)};
    \node[font=\small\bfseries, text=qaaccent2] (h2) at (7,0) {Transaction T2 (client 2)};
    \node[db] (hb) at (3.5,0) {stock en base};
    \node[tx]  (a1) at (0,-0.8) {SELECT stock $\to$ 10};
    \node[db]  (b1) at (3.5,-0.8) {10};
    \node[txb] (c1) at (7,-1.5) {SELECT stock $\to$ 10};
    \node[db]  (b2) at (3.5,-1.5) {10};
    \node[txb] (c2) at (7,-2.2) {UPDATE stock = 10-3};
    \node[txb] (c3) at (7,-2.9) {COMMIT};
    \node[db]  (b3) at (3.5,-2.9) {7};
    \node[tx]  (a2) at (0,-3.6) {UPDATE stock = 10-2};
    \node[tx]  (a3) at (0,-4.3) {COMMIT};
    \node[db, draw=qaaccent2, text=qaaccent2]  (b4) at (3.5,-4.3) {\textbf{8} (au lieu de 5)};
    \draw[timeline] (-2.4,-0.4) -- node[left, font=\scriptsize, text=qamuted, rotate=90, anchor=south] {temps} (-2.4,-4.6);
  \end{tikzpicture}
  \end{center}
  \begin{itemize}
    \item Chaque transaction calcule le nouveau stock \textbf{dans le code Java} \`a partir d'une valeur lue trop t\^ot
    \item C'est le niveau d'isolation \textbf{par d\'efaut} de PostgreSQL (\texttt{READ COMMITTED}) : \alert{test\'e, r\'esultat 8}
  \end{itemize}
\end{frame}

\begin{frame}{Les anomalies de concurrence}
  \begin{datacard}
    \small
    \begin{tabular}{@{}p{3.6cm}p{8.8cm}@{}}
      \toprule
      \textbf{Anomalie} & \textbf{Ce qui se passe} \\
      \midrule
      Lecture sale\newline{\scriptsize\color{qamuted}\textit{dirty read} (P1)} & T1 lit une valeur modifi\'ee par T2 \textbf{non encore valid\'ee} ; T2 fait ensuite \texttt{ROLLBACK} \\
      Lecture non r\'ep\'etable\newline{\scriptsize\color{qamuted}\textit{non-repeatable read} (P2)} & T1 lit deux fois la m\^eme ligne et obtient deux valeurs diff\'erentes (T2 a valid\'e entre-temps) \\
      Lecture fant\^ome\newline{\scriptsize\color{qamuted}\textit{phantom read} (P3)} & T1 relance une requ\^ete et voit des \textbf{lignes en plus ou en moins} (insertion/suppression de T2) \\
      Mise \`a jour perdue\newline{\scriptsize\color{qamuted}\textit{lost update} (P4)} & T1 et T2 lisent puis \'ecrivent la m\^eme ligne : l'\'ecriture de l'une \'ecrase celle de l'autre \\
      \bottomrule
    \end{tabular}
  \end{datacard}
  \vspace{1mm}
  \source{P1 \`a P3 : ph\'enom\`enes de la norme ANSI SQL-92. P4 : ajout\'e par Berenson et al. (SIGMOD 1995), absent de la norme.}
\end{frame}

\begin{frame}{Exemple : une lecture incoh\'erente (\textit{read skew})}
  \begin{center}
  \begin{tikzpicture}
    \node[font=\small\bfseries, text=qaaccent] at (0,0) {T1 : transfert de 30 (A $\to$ B)};
    \node[font=\small\bfseries, text=qaaccent2] at (7,0) {T2 : affiche A + B};
    \node[db] at (3.5,0) {total en base : 100};
    \node[tx]  at (0,-0.8) {UPDATE A = A - 30};
    \node[db]  at (3.5,-0.8) {A = 60, B = 40};
    \node[txb] at (7,-1.5) {SELECT A $\to$ 60};
    \node[tx]  at (0,-2.2) {UPDATE B = B + 30};
    \node[tx]  at (0,-2.9) {COMMIT};
    \node[db]  at (3.5,-2.9) {A = 30, B = 70};
    \node[txb] at (7,-3.6) {SELECT B $\to$ 70};
    \node[db, draw=qaaccent2, text=qaaccent2] at (3.5,-4.3) {\textbf{130} (au lieu de 100)};
    \node[font=\scriptsize, text=qaaccent2, anchor=west] at (5.2,-4.3) {$\leftarrow$ 60 + 70};
    \draw[timeline] (-2.4,-0.4) -- node[left, font=\scriptsize, text=qamuted, rotate=90, anchor=south] {temps} (-2.4,-4.6);
  \end{tikzpicture}
  \end{center}
  \begin{itemize}
    \item A et B : deux entrep\^ots. PostgreSQL ne lit \textbf{jamais} une valeur non valid\'ee, mais deux \texttt{SELECT} s\'epar\'es peuvent voir \textbf{deux \'etats} de la base
    \item Rem\`ede : \textbf{une seule requ\^ete} (\texttt{SELECT SUM(stock) \ldots{}}) ou \texttt{REPEATABLE READ}
  \end{itemize}
\end{frame}

% ================================================================
\section{Niveaux d'isolation}

\begin{frame}{Les quatre niveaux d'isolation}
  \begin{datacard}
    \small
    \begin{tabular}{@{}lcccc@{}}
      \toprule
      \textbf{Niveau} & \textbf{Lect. sale} & \textbf{Non r\'ep\'etable} & \textbf{Fant\^ome} & \textbf{M\`aj perdue} \\
      \midrule
      \texttt{READ UNCOMMITTED} & possible$^{*}$ & possible & possible & possible \\
      \texttt{READ COMMITTED} (d\'efaut PG) & --- & possible & possible & \alert{possible} \\
      \texttt{REPEATABLE READ} & --- & --- & possible$^{*}$ & d\'etect\'ee \\
      \texttt{SERIALIZABLE} & --- & --- & --- & d\'etect\'ee \\
      \bottomrule
    \end{tabular}
  \end{datacard}
  \vspace{1mm}
  {\footnotesize\color{qamuted} $^{*}$ selon la norme. PostgreSQL est plus strict : \texttt{READ UNCOMMITTED} s'y comporte comme \texttt{READ COMMITTED}
  (jamais de lecture sale), et \texttt{REPEATABLE READ} y emp\^eche aussi les fant\^omes.
  La colonne \og M\`aj perdue \fg{} n'est pas dans la norme (P4, Berenson et al. 1995).
  Source : documentation PostgreSQL, \S13.2.}
  \begin{itemize}
    \item Plus le niveau est \'elev\'e, plus il y a de transactions \textbf{refus\'ees} ou \textbf{bloqu\'ees} : on paie l'isolation en d\'ebit
    \item PostgreSQL utilise le \textbf{MVCC} : chaque transaction lit une \og photo \fg{} (snapshot) de la base ; les lecteurs ne bloquent jamais les \'ecrivains
  \end{itemize}
\end{frame}

\begin{frame}[fragile]{Choisir le niveau d'isolation}
  \begin{columns}[T]
    \begin{column}{0.48\textwidth}
      \textbf{En SQL}
\begin{lstlisting}[language=SQL]
BEGIN ISOLATION LEVEL REPEATABLE READ;
SELECT stock FROM Produit
 WHERE numProduit = 1;
UPDATE Produit SET stock = 8
 WHERE numProduit = 1;
COMMIT;
\end{lstlisting}
    \end{column}
    \begin{column}{0.48\textwidth}
      \textbf{En JDBC}
\begin{lstlisting}[language=Java]
c.setAutoCommit(false);
c.setTransactionIsolation(
  Connection.TRANSACTION_REPEATABLE_READ);
// ... requetes ...
c.commit();
\end{lstlisting}
    \end{column}
  \end{columns}
  \vspace{2mm}
  \begin{block}{Ce qu'on a observ\'e sur PostgreSQL (sc\'enario des deux clients)}
    \begin{itemize}
      \item \texttt{READ COMMITTED} : stock final 8 (au lieu de 5) \hfill \alert{mise \`a jour perdue}
      \item \texttt{REPEATABLE READ} : T1 re\c{c}oit \texttt{ERROR: could not serialize access due to concurrent update} \hfill SQLState \texttt{40001}
    \end{itemize}
    $\Rightarrow$ \`a un niveau \'elev\'e, l'application doit \textbf{rejouer} la transaction refus\'ee
  \end{block}
\end{frame}

% ================================================================
\section{\'Eviter les mises \`a jour perdues}

\begin{frame}[fragile]{Solution 1 : laisser la base faire le calcul}
\begin{lstlisting}[language=Java]
// NON : lecture, calcul en Java, puis ecriture (2 requetes, fenetre de concurrence)
int stock = lireStock(c, num);
majStock(c, num, stock - quantite);

// OUI : une seule requete, atomique pour la ligne concernee
PreparedStatement ps = c.prepareStatement(
  "UPDATE Produit SET stock = stock - ? WHERE numProduit = ? AND stock >= ?");
ps.setInt(1, quantite); ps.setInt(2, num); ps.setInt(3, quantite);
if (ps.executeUpdate() == 0) throw new StockInsuffisantException(num);
\end{lstlisting}
  \begin{itemize}
    \item La 2\textsuperscript{e} transaction attend le verrou de ligne, puis relit la valeur \`a jour : r\'esultat 5
    \item La condition \texttt{stock >= ?} v\'erifie la r\`egle m\'etier \textbf{dans la m\^eme requ\^ete}
    \item \`A privil\'egier d\`es que le calcul s'exprime en SQL
  \end{itemize}
\end{frame}

\begin{frame}{Les verrous : partag\'e (S) et exclusif (X)}
  \begin{columns}[T]
    \begin{column}{0.40\textwidth}
      \begin{exampleblock}{Compatibilit\'e}
        \centering\small
        \begin{tabular}{@{}l|c|c@{}}
          & \textbf{d\'ej\`a S} & \textbf{d\'ej\`a X} \\ \hline
          \textbf{veut S} & \alert{ok} & attend \\
          \textbf{veut X} & attend & attend
        \end{tabular}
      \end{exampleblock}
    \end{column}
    \begin{column}{0.56\textwidth}
      \begin{exampleblock}{Dans PostgreSQL (par ligne)}
        \small
        \texttt{SELECT} simple : aucun verrou\\[0.4ex]
        \texttt{FOR SHARE} : verrou S\\[0.4ex]
        \texttt{UPDATE}, \texttt{DELETE}, \texttt{FOR UPDATE} : verrou X
      \end{exampleblock}
    \end{column}
  \end{columns}
  \vspace{2mm}
  \begin{itemize}
    \item Les lecteurs simples ne bloquent jamais les \'ecrivains (MVCC)
    \item Verrouillage \`a deux phases \textbf{strict} : les verrous sont gard\'es jusqu'au \texttt{COMMIT} ou \texttt{ROLLBACK}
    \item D'o\`u : transactions \textbf{courtes}, et toujours le \textbf{m\^eme ordre} de verrouillage
  \end{itemize}
  \vfill
  \source{Source : documentation PostgreSQL, \S13.3 \og Explicit Locking \fg{} (verrous de ligne, interblocages).}
\end{frame}

\begin{frame}[fragile]{Solution 2 : verrou pessimiste (\texttt{SELECT \ldots{} FOR UPDATE})}
\begin{lstlisting}[language=Java]
c.setAutoCommit(false);
PreparedStatement sel = c.prepareStatement(
  "SELECT stock FROM Produit WHERE numProduit = ? FOR UPDATE");  // pose un verrou
sel.setInt(1, num);
ResultSet rs = sel.executeQuery(); rs.next();
int stock = rs.getInt("stock");
// ... regles metier complexes en Java ...
majStock(c, num, stock - quantite);
c.commit();                                                    // libere le verrou
\end{lstlisting}
  \begin{itemize}
    \item Les autres transactions qui veulent la m\^eme ligne \textbf{attendent} le \texttt{COMMIT}
    \item Utile quand la logique est trop complexe pour une seule requ\^ete
    \item Transaction \textbf{courte} obligatoire : jamais d'attente utilisateur pendant qu'on tient un verrou
  \end{itemize}
\end{frame}

\begin{frame}{Avec \texttt{FOR UPDATE}, la 2\textsuperscript{e} transaction attend}
  \begin{center}
  \begin{tikzpicture}
    \node[font=\small\bfseries, text=qaaccent] at (0,0) {T1 : commande 2 articles};
    \node[font=\small\bfseries, text=qaaccent2] at (7,0) {T2 : commande 3 articles};
    \node[db] at (3.5,0) {stock en base};
    \node[tx]  at (0,-0.8) {SELECT \ldots{} FOR UPDATE $\to$ 10};
    \node[db]  at (3.5,-0.8) {10};
    \node[txb, dashed] at (7,-1.5) {SELECT \ldots{} FOR UPDATE};
    \node[font=\scriptsize, text=qaaccent2] at (7,-1.95) {attend le verrou de T1};
    \node[tx]  at (0,-2.5) {UPDATE stock = 10-2};
    \node[tx]  at (0,-3.2) {COMMIT};
    \node[db]  at (3.5,-3.2) {8};
    \node[txb] at (7,-3.2) {verrou obtenu $\to$ 8};
    \node[txb] at (7,-3.9) {UPDATE stock = 8-3};
    \node[txb] at (7,-4.6) {COMMIT};
    \node[db, draw=qaaccent2, text=qaaccent2] at (3.5,-4.6) {\textbf{5} (correct)};
    \draw[timeline] (-2.4,-0.4) -- node[left, font=\scriptsize, text=qamuted, rotate=90, anchor=south] {temps} (-2.4,-4.9);
  \end{tikzpicture}
  \end{center}
  \begin{itemize}\small
    \item \texttt{READ COMMITTED} (d\'efaut) : T2 est \textbf{bloqu\'ee} jusqu'au \texttt{COMMIT} de T1, puis lit la valeur \textbf{\`a jour}
    \item Prix \`a payer : les deux transactions s'ex\'ecutent l'une apr\`es l'autre
    \item \texttt{REPEATABLE READ} : T2 recevrait plut\^ot l'erreur \texttt{40001}, \`a rejouer
  \end{itemize}
\end{frame}

\begin{frame}[fragile]{Solution 3 : verrou optimiste (colonne de version)}
\begin{lstlisting}[language=SQL]
ALTER TABLE Produit ADD COLUMN version INT NOT NULL DEFAULT 0;
-- lecture : on memorise la version lue (ex. 7)
SELECT stock, version FROM Produit WHERE numProduit = 1;
-- ecriture : seulement si personne n'a modifie la ligne entre-temps
UPDATE Produit SET stock = 8, version = version + 1
 WHERE numProduit = 1 AND version = 7;
\end{lstlisting}
  \begin{itemize}
    \item 0 ligne modifi\'ee $\Rightarrow$ conflit : on pr\'evient l'utilisateur ou on rejoue
    \item Aucun verrou tenu entre la lecture et l'\'ecriture : adapt\'e aux \textbf{longues interactions} (formulaire web)
    \item JPA l'automatise avec \texttt{@Version} (\alert{chapitre 6})
  \end{itemize}
\end{frame}

\begin{frame}{Quelle solution choisir ?}
  \begin{datacard}
    \small
    \begin{tabular}{@{}p{3.6cm}p{4.4cm}p{4.4cm}@{}}
      \toprule
      \textbf{Approche} & \textbf{Quand l'utiliser} & \textbf{Prix \`a payer} \\
      \midrule
      Calcul en SQL & mise \`a jour simple (compteur, stock) & aucun \\
      \texttt{FOR UPDATE} & logique complexe, conflits fr\'equents & attente, risque d'interblocage \\
      Version (optimiste) & conflits rares, interaction longue & g\'erer l'\'echec (message, rejeu) \\
      \texttt{REPEATABLE READ}/ \texttt{SERIALIZABLE} & r\`egles portant sur plusieurs lignes & rejouer sur SQLState \texttt{40001} \\
      \bottomrule
    \end{tabular}
  \end{datacard}
\end{frame}

% ================================================================
\section{Interblocages}

\begin{frame}{Interblocage (\textit{deadlock})}
  \begin{center}
  \begin{tikzpicture}
    \node[tx]  (a1) at (0,0) {UPDATE produit 1};
    \node[txb] (b1) at (6,-0.6) {UPDATE produit 2};
    \node[tx]  (a2) at (0,-1.2) {UPDATE produit 2 \ldots{} attend T2};
    \node[txb] (b2) at (6,-1.8) {UPDATE produit 1 \ldots{} attend T1};
    \node[font=\small\bfseries, text=qaaccent] at (0,0.6) {T1};
    \node[font=\small\bfseries, text=qaaccent2] at (6,0.6) {T2};
    \draw[-Latex, thick, qaaccent2] (b2.west) -- (a1.east);
    \draw[-Latex, thick, qaaccent] (a2.east) -- (b1.west);
  \end{tikzpicture}
  \end{center}
  \begin{itemize}
    \item PostgreSQL d\'etecte le cycle (au bout d'une seconde) et \textbf{annule l'une des deux} :
      \texttt{ERROR: deadlock detected} \hfill SQLState \texttt{40P01}
    \item Pr\'evention : toujours verrouiller dans le \textbf{m\^eme ordre} (ex. produit de plus petit num\'ero d'abord), transactions courtes
    \item L'application traite \texttt{40P01} comme \texttt{40001} : on \textbf{rejoue} la transaction
  \end{itemize}
\end{frame}

\begin{frame}[fragile]{Rejouer une transaction refus\'ee}
\begin{lstlisting}[language=Java]
for (int essai = 1; ; essai++) {
  try (Connection c = dataSource.getConnection()) {
    c.setAutoCommit(false);
    c.setTransactionIsolation(Connection.TRANSACTION_REPEATABLE_READ);
    try {
      passerCommande(c, client, panier);
      c.commit();
      return;
    } catch (SQLException e) {
      c.rollback();
      String etat = e.getSQLState();
      boolean rejouable = "40001".equals(etat) || "40P01".equals(etat);
      if (!rejouable || essai == 3) throw e;
    }
  }
}
\end{lstlisting}
\end{frame}

% ================================================================
\section{Contraintes, erreurs et bonnes pratiques}

\begin{frame}[fragile]{Contraintes d\'eclaratives}
\begin{lstlisting}[language=SQL,basicstyle=\lstcode\scriptsize]
ALTER TABLE Etudiant ALTER COLUMN nom SET NOT NULL;
ALTER TABLE Etudiant ADD CONSTRAINT chk_moy CHECK (moyenne BETWEEN 0 AND 20);
ALTER TABLE Etudiant ADD CONSTRAINT uq_mail UNIQUE (email);
ALTER TABLE Inscription ALTER COLUMN statut SET DEFAULT 'en attente';
\end{lstlisting}
  \begin{itemize}
    \item \texttt{NOT NULL}, \texttt{CHECK}, \texttt{UNIQUE}, cl\'es primaires et \'etrang\`eres : la \textbf{base refuse} toute ligne qui les viole
    \item \texttt{DEFAULT} n'interdit rien : il fournit la valeur utilis\'ee quand la colonne est omise \`a l'\texttt{INSERT}
    \item \texttt{ALTER TABLE} modifie le sch\'ema sans le recr\'eer ; ajouter un \texttt{NOT NULL}, un \texttt{CHECK} ou un \texttt{UNIQUE} \'echoue si des lignes existantes violent d\'ej\`a la r\`egle
    \item Relation plusieurs-\`a-plusieurs : \textbf{table d'association} (ex. \texttt{Inscription(numEtudiant, numCours)}) \`a cl\'e primaire composite
  \end{itemize}
\end{frame}

\begin{frame}[fragile]{Erreurs SQL : lire le \texttt{SQLState}}
  \begin{datacard}
    \footnotesize
    \begin{tabular}{@{}p{1.8cm}p{6.2cm}p{2.6cm}@{}}
      \toprule
      \textbf{SQLState} & \textbf{Signification} & \textbf{R\'eaction} \\
      \midrule
      \texttt{23505} & doublon (cl\'e primaire, \texttt{UNIQUE}) & erreur m\'etier \\
      \texttt{23503} & cl\'e \'etrang\`ere inexistante & erreur m\'etier \\
      \texttt{23514} & \texttt{CHECK} viol\'e (ex. \texttt{moyenne > 20}) & erreur m\'etier \\
      \texttt{40001} & \'echec de s\'erialisation & \textbf{rejouer} \\
      \texttt{40P01} & interblocage d\'etect\'e & \textbf{rejouer} \\
      \bottomrule
    \end{tabular}
  \end{datacard}
\begin{lstlisting}[language=Java,basicstyle=\lstcode\scriptsize]
catch (SQLException e) {
  c.rollback();
  if ("23505".equals(e.getSQLState()))
    throw new EtudiantExistantException(num);
  throw e;
}
\end{lstlisting}
  \begin{itemize}
    \item \texttt{PSQLException} n'a pas de sous-classe d\'edi\'ee : on lit \texttt{getSQLState()}
  \end{itemize}
  \source{Source : documentation PostgreSQL, annexe A \og PostgreSQL Error Codes \fg{} et \S13.5.}
\end{frame}

\begin{frame}{Bonnes pratiques}
  \begin{itemize}
    \item Des transactions \textbf{courtes} : pas d'appel r\'eseau, pas d'attente utilisateur \`a l'int\'erieur
    \item La fronti\`ere de transaction est dans la \textbf{couche service} (chapitre 5), pas dans chaque DAO
    \item Pr\'ef\'erer les mises \`a jour \textbf{relatives} en SQL (\texttt{stock = stock - ?})
    \item Faire v\'erifier les r\`egles par la base : \texttt{CHECK (stock >= 0)}, \texttt{UNIQUE}, cl\'es \'etrang\`eres
    \item Montants : \texttt{NUMERIC} en SQL, \texttt{BigDecimal} en Java : \texttt{REAL}/\texttt{double} sont inexacts ($0{,}1 + 0{,}2 \neq 0{,}3$)
    \item Toujours \texttt{rollback()} en cas d'exception, et \texttt{close()} pour rendre la connexion au pool
  \end{itemize}
\end{frame}

% ================================================================
\section{Exigences pour les rendus de TP}

\begin{frame}[fragile]{Secrets : le fichier \texttt{.env}}
  \begin{columns}[T]
    \begin{column}{0.48\textwidth}
      \textbf{.env} (\alert{jamais dans Git})
\begin{lstlisting}[language={},basicstyle=\lstcode\scriptsize]
POSTGRES_USER=app
POSTGRES_PASSWORD=change-me
POSTGRES_DB=campus_db
DB_URL=jdbc:postgresql://db:5432/campus_db
DB_USER=app
DB_PASSWORD=change-me
\end{lstlisting}
    \end{column}
    \begin{column}{0.48\textwidth}
      \textbf{Lecture c\^ot\'e Java}
\begin{lstlisting}[language=Java,basicstyle=\lstcode\scriptsize]
// conteneur : injecte par Docker
System.getenv("DB_URL");
// hors Docker : dotenv-java
Dotenv.configure()
  .ignoreIfMissing().load()
  .get("DB_URL");
\end{lstlisting}
    \end{column}
  \end{columns}
  \vspace{2mm}
  \begin{itemize}
    \item \texttt{.env} dans \texttt{.gitignore} ; on versionne un \texttt{.env.example} sans secret
    \item Docker Compose injecte \texttt{.env} dans le conteneur (\texttt{env\_file})
    \item Hors Docker (IntelliJ) : l'h\^ote \texttt{db} devient \texttt{localhost}
  \end{itemize}
\end{frame}

\begin{frame}{Un rendu de qualit\'e, c'est\ldots}
  \begin{columns}[T]
    \begin{column}{0.48\textwidth}
      \begin{block}{D\'ep\^ot Git}
        \begin{itemize}
          \item Un d\'ep\^ot par bin\^ome, commits r\'eguliers et messages explicites
          \item \texttt{.gitignore} : \texttt{target/}, \texttt{.idea/}, \alert{\texttt{.env}}
          \item Aucun mot de passe dans le code ni dans l'historique
          \item Projet Maven standard (\texttt{src/main/java}, \texttt{src/test/java})
        \end{itemize}
      \end{block}
    \end{column}
    \begin{column}{0.48\textwidth}
      \begin{block}{README.md}
        \begin{itemize}
          \item Comment lancer la base (\texttt{cp .env.example .env}, puis \texttt{docker compose up -d})
          \item Comment compiler et ex\'ecuter (\texttt{mvn \ldots})
          \item Ce qui est fait, ce qui ne l'est pas
          \item Choix techniques et difficult\'es
        \end{itemize}
      \end{block}
    \end{column}
  \end{columns}
  \vspace{2mm}
  \begin{alertblock}{Test d'acceptation}
    Un enseignant clone le d\'ep\^ot sur une machine vierge et suit le README : \textbf{\c{c}a doit fonctionner}.
  \end{alertblock}
\end{frame}

\begin{frame}{R\'ef\'erences (1/2) : articles et livre}
  \footnotesize
  \begin{itemize}\setlength{\itemsep}{1.2ex}
    \item H. Berenson, P. Bernstein, J. Gray, J. Melton, E. O'Neil, P. O'Neil.
      \textit{A Critique of ANSI SQL Isolation Levels}. ACM SIGMOD 1995, pp.~1--10.
      \url{https://doi.org/10.1145/223784.223785}\\
      {\color{qamuted}Phénom\`enes P1--P4, \textit{read skew}, \textit{snapshot isolation}.}
    \item T. H\"arder, A. Reuter.
      \textit{Principles of Transaction-Oriented Database Recovery}. ACM Computing Surveys 15(4), 1983, pp.~287--317.
      \url{https://doi.org/10.1145/289.291}\\
      {\color{qamuted}R\'ef\'erence classique sur les transactions et la reprise apr\`es panne.}
    \item M. Kleppmann. \textit{Designing Data-Intensive Applications}. O'Reilly, 2017.
      \url{https://dataintensive.net/}\\
      {\color{qamuted}Chapitre 7 \og Transactions \fg{} : isolation, mises \`a jour perdues, \textit{write skew}.}
  \end{itemize}
\end{frame}

\begin{frame}{R\'ef\'erences (2/2) : documentation et outils}
  \footnotesize
  \begin{itemize}\setlength{\itemsep}{1.0ex}
    \item PostgreSQL, chapitre 13 \og Concurrency Control \fg{} : \S13.2 \textit{Transaction Isolation},
      \S13.3 \textit{Explicit Locking}, \S13.5 \textit{Serialization Failure Handling}\\
      \url{https://www.postgresql.org/docs/current/transaction-iso.html}\\
      \url{https://www.postgresql.org/docs/current/explicit-locking.html}\\
      \url{https://www.postgresql.org/docs/current/mvcc-serialization-failure-handling.html}
    \item PostgreSQL, annexe A \og PostgreSQL Error Codes \fg{} (SQLState)\\
      \url{https://www.postgresql.org/docs/current/errcodes-appendix.html}
    \item Oracle, \textit{The Java Tutorials} : \og Using Transactions \fg{} (JDBC)\\
      \url{https://docs.oracle.com/javase/tutorial/jdbc/basics/transactions.html}
    \item Pilote JDBC PostgreSQL : classe \texttt{PSQLException}\\
      \url{https://jdbc.postgresql.org/documentation/publicapi/org/postgresql/util/PSQLException.html}
    \item Outils : HikariCP \url{https://github.com/brettwooldridge/HikariCP} ;
      dotenv-java \url{https://github.com/cdimascio/dotenv-java} ;
      Docker Compose, \texttt{env\_file} \url{https://docs.docker.com/compose/how-tos/environment-variables/set-environment-variables/}
  \end{itemize}
\end{frame}

\rquestionslide

\end{document}
